\documentclass[11pt,a4paper]{article}
\usepackage[T1]{fontenc}
\usepackage[utf8]{inputenc}
\usepackage{lmodern,microtype}
\usepackage[margin=27mm,headheight=14pt]{geometry}
\usepackage{amsmath,amssymb,amsthm,mathtools}
\usepackage{booktabs,array,longtable}
\usepackage{xcolor,xurl,hyperref,fancyhdr}
\usepackage{enumitem}
\hypersetup{colorlinks=true,linkcolor=blue!45!black,citecolor=blue!45!black,urlcolor=blue!45!black,
 pdftitle={Fixed-composition excursions: D-finiteness, transcendence, and all-order asymptotics for OEIS A215561},
 pdfauthor={Research report prepared for Vladimir Reshetnikov}}
\pagestyle{fancy}\fancyhf{}
\fancyhead[L]{\small Fixed-composition excursions}
\fancyhead[R]{\small OEIS A215561}
\fancyfoot[C]{\thepage}
\setlength{\parskip}{4pt}
\setlength{\parindent}{1.2em}
\setlength{\emergencystretch}{2em}
\numberwithin{equation}{section}
\newtheorem{theorem}{Theorem}[section]
\newtheorem{proposition}[theorem]{Proposition}
\newtheorem{lemma}[theorem]{Lemma}
\newtheorem{corollary}[theorem]{Corollary}
\theoremstyle{definition}\newtheorem{definition}[theorem]{Definition}
\theoremstyle{remark}\newtheorem{remark}[theorem]{Remark}
\newcommand{\N}{\mathbb N}
\newcommand{\Q}{\mathbb Q}
\newcommand{\R}{\mathbb R}
\newcommand{\C}{\mathbb C}
\newcommand{\A}{\overline{\mathbb Q}}
\newcommand{\ct}{\operatorname{CT}}
\newcommand{\diag}{\operatorname{Diag}}
\newcommand{\E}{\mathcal E}
\newcommand{\Lcal}{\mathcal L}
\newcommand{\Prob}{\mathbb P}
\newcommand{\Exp}{\mathbb E}
\newcommand{\Var}{\operatorname{Var}}
\newcommand{\Cov}{\operatorname{Cov}}
\newcommand{\coeff}[2]{[#1]#2}
\newcommand{\dd}{\,\mathrm d}
\newcommand{\one}{\boldsymbol 1}
\DeclareMathOperator{\Rea}{Re}
\DeclareMathOperator{\Res}{Res}
\begin{document}
\begin{titlepage}
\centering
\vspace*{14mm}
{\large OEIS RESEARCH REPORT\par}
\vspace{11mm}
{\LARGE\bfseries Fixed-composition excursions\par}
\vspace{5mm}
{\Large D-finiteness, transcendence, and all-order\par asymptotics for OEIS A215561\par}
\vspace{12mm}
{\large With an explicit correction for A215570\par and asymptotic inversion\par}
\vspace{14mm}
{Prepared for Vladimir Reshetnikov\par}
{October 1, 2026\par}
\vspace{10mm}
\begin{minipage}{0.89\textwidth}
\begin{abstract}
Let $A_r(n)$ count words containing $n$ copies of each letter in
$\{1,\ldots,r\}$ whose every prefix has sum at most its average over all
such words. We prove the fixed-row asymptotic conjecture recorded in
OEIS A215561, identify its constant as a positive algebraic number, and
establish a complete expansion in inverse powers of $n$. The central
formula is
\[
 A_r(n)\sim \frac{\kappa_r}{rn}\frac{(rn)!}{(n!)^r},
 \qquad
 \kappa_r=\exp\!\left(\sum_{m\ge1}\frac{\Prob(S_m=0)}m\right),
\]
where $S_m$ is the centered, uniformly weighted walk associated with the
alphabet. A finite kernel-root product computes $\kappa_r$.

We also prove that every fixed row is D-finite, while its ordinary
generating function is transcendental for every $r\ge3$. This separates
recurrence existence from algebraicity and from the still-unproved
specific recurrence conjecture for A215570. An explicit first correction
for that sequence, a fixed-composition contact law, and an all-order
Lambert-$W$ inversion procedure are derived. Exact enumeration and
symbolic checks accompany the proofs.
\end{abstract}
\end{minipage}
\vfill
\begin{minipage}{0.89\textwidth}\small
\textbf{Status.} This is a newly written mathematical research manuscript,
not a peer-reviewed publication or a Lean-verified development. Classical
kernel and diagonal theorems are credited separately from the applications
and deductions proved here. No claim of exhaustive bibliographic priority
is made. The particular order-three recurrence conjecture of Kauers and
Koutschan is \emph{not} claimed to be proved.
\end{minipage}
\end{titlepage}
\tableofcontents
\newpage

\section{The problem, its significance, and the source audit}
\subsection{A natural multiset-permutation family}
For a fixed integer $r\ge2$, let $A_r(n)$ be the number of words
$w=w_1\cdots w_{rn}$ containing exactly $n$ occurrences of each letter
$1,\ldots,r$ and satisfying
\begin{equation}\label{eq:prefix}
 \sum_{j=1}^k w_j\le \frac{r+1}{2}k\qquad(0\le k\le rn).
\end{equation}
These are rows of the array A215561. The average on the right is exact:
in a uniformly chosen multiset permutation, each position has expected
letter $(r+1)/2$. This family includes Catalan numbers and a balanced
Motzkin-word enumeration, but already the fourth and fifth rows lead to
substantially more difficult coefficient extraction.

The relevant sequence identifiers are
\begin{center}
\begin{tabular}{ccl}
\toprule
$r$ & OEIS entry & Description of the row \\
\midrule
2&A000108&Catalan numbers\\
3&A007004&Balanced three-letter excursions\\
4&A215562&Four-letter prefix-average words\\
5&A215570&Five-letter prefix-average words\\
6&A215571&Six-letter prefix-average words\\
7&A215593&Seven-letter prefix-average words\\
\bottomrule
\end{tabular}
\end{center}
The array entry records Kotesovec's 2016 conjecture that, for each fixed
$r>1$,
\begin{equation}\label{eq:oeisconj}
 A_r(n)\sim c(r)\frac{r^{rn}}{\pi^{(r-1)/2}n^{(r+1)/2}}.
\end{equation}
The fourth and fifth entries already supply explicit leading constants;
those constants are credited to Kotesovec, rather than presented here as
new discoveries \cite{oeisarray,oeis4,oeis5}.

\subsection{What is and is not resolved}
Kauers and Koutschan discuss A215570 and A215562 in Section 6.1 of their
2023 paper. Their finite transfer matrices grow with $n$, which prevents
their argument from proving D-finiteness. They leave a particular
order-three recurrence for A215570 conjectural and ask for a provably
correct recurrence for A215562 \cite{kk}. We replace the growing matrix
by a fixed algebraic generating function in composition-marking variables.

The main results below establish the existence of polynomial-coefficient
recurrences for \emph{every} fixed row. They do not print a recurrence for
the fourth row and do not certify the particular guessed operator for
the fifth. Recurrence existence, an explicit annihilating operator, and
minimality of that operator are three different assertions.

The source check used the live OEIS entries, the cited 2023 paper, the
classical directed-path kernel theory, and targeted searches for the
sequence identifiers. The fixed-row asymptotic was still labelled a
conjecture in A215561 when checked. No proof of the precise results below
was identified in that search. This is a bounded source audit, not a
proof that no such result exists anywhere in the literature.

\subsection{Relationship to ProveIt}
The inspected ProveIt snapshot was
\begin{center}\small\texttt{40b17f2fc79a10a6fc1a50ca9c0c3d6f77be69e5}.\end{center}
Its transseries-and-inversion documentation distinguishes exact cores,
finite-order expansions, remainder transport, and discrete staircase
inversion. We follow that organization in Section~\ref{sec:inverse}.
The current path inspected was
\begin{quote}\small\path{Analysis/Transseries/docs/series-and-transseries/Transseries_And_Inversion/README.md}.
\end{quote}
The document's own formalization inventory is deliberately narrower than
its mathematical exposition \cite{proveit}. No repository theorem is
silently imported as a formal proof of our excursion results. Searches
for A215561 and A215562 returned no repository matches; that limited check
is not an audit of every file or of other identifiers for the same problem.

\section{Main theorems}
Put
\begin{equation}\label{eq:steps}
 g_r=\gcd(2,r-1),\quad
 s_i=\frac{2i-r-1}{g_r},\quad
 h=-s_1,\quad
 P_r(u)=\frac1r\sum_{i=1}^r u^{s_i}.
\end{equation}
For odd $r=2h+1$, the steps are $-h,-h+1,\ldots,h$. For even $r$, they
are the odd integers $-(r-1),-(r-3),\ldots,r-1$. Let $S_m$ be the sum of
$m$ independent steps chosen uniformly from this set.

\begin{theorem}[All fixed rows: leading term and complete expansion]\label{thm:main}
For every fixed $r\ge2$, the series
\begin{equation}\label{eq:kappaseries}
 \log\kappa_r=\sum_{m\ge1}\frac{\Prob(S_m=0)}m
\end{equation}
converges, and $\kappa_r$ is a positive real algebraic number. For each
integer $M\ge0$ there are real algebraic numbers $\alpha_{r,1},\ldots,
\alpha_{r,M}$ such that
\begin{equation}\label{eq:full}
 A_r(n)=K_r\frac{r^{rn}}{n^{(r+1)/2}}
 \left(1+\sum_{j=1}^M\frac{\alpha_{r,j}}{n^j}
       +O_{r,M}(n^{-M-1})\right),
\end{equation}
where
\begin{equation}\label{eq:constant}
 K_r=\frac{\kappa_r}{\sqrt r\,(2\pi)^{(r-1)/2}},\qquad
 c(r)=\frac{\kappa_r}{\sqrt r\,2^{(r-1)/2}}.
\end{equation}
In particular,
\begin{equation}\label{eq:probability}
 \frac{A_r(n)}{(rn)!/(n!)^r}\sim\frac{\kappa_r}{rn}.
\end{equation}
All coefficients in \eqref{eq:full} are computable by finite algebraic
operations on kernel expansions and Gaussian moments.
\end{theorem}
The constants implied by the $O$-notation are allowed to depend on $r$
and $M$. Nothing in this statement is uniform when the alphabet size
grows with $n$.

\begin{theorem}[Algebraic ambient series; transcendental row series]\label{thm:class}
The multivariate generating function counting nonnegative excursions
with steps \eqref{eq:steps}, marking every step type independently, is
algebraic over $\Q(x_1,\ldots,x_r)$. Consequently
\[
 F_r(z)=\sum_{n\ge0}A_r(n)z^n
\]
is D-finite for every fixed $r$. The function $F_r$ is transcendental over
$\Q(z)$ for every $r\ge3$. For $r=2$ it is the algebraic Catalan generating
function.
\end{theorem}
Thus taking a diagonal preserves D-finiteness here, but does not preserve
algebraicity. The transcendence assertion uses not just the growth
exponent but, for even $r$, the power of $\pi$ in the leading constant.

\begin{theorem}[An explicit next term for A215570]\label{thm:five}
For the fifth row,
\begin{equation}\label{eq:fivecorrection}
 A_5(n)=\frac{3\sqrt5-5}{8\pi^2}\frac{5^{5n}}{n^3}
 \left(1-\frac{13(5-\sqrt5)}{50n}+O(n^{-2})\right).
\end{equation}
Moreover every coefficient $\alpha_{5,j}$ in its full expansion belongs
to $\Q(\sqrt5)$.
\end{theorem}
The leading factor in \eqref{eq:fivecorrection} is the previously posted
OEIS factor. The proof and the correction are developed below without
assuming the conjectured recurrence.

\section{The exact excursion generating function}
\subsection{Encoding the prefix condition}
Replacing a letter $i$ by $r+1-i$ reverses the prefix inequality and preserves
the multiset. It is therefore equivalent to count words whose centered
partial sums, using steps $s_i$, are \emph{nonnegative}. The total centered
sum is zero. These are one-dimensional excursions with prescribed numbers
of the individual steps.

For a count vector $\boldsymbol k=(k_1,\ldots,k_r)$ write
$|\boldsymbol k|=\sum_i k_i$ and
$\boldsymbol x^{\boldsymbol k}=\prod_i x_i^{k_i}$.
Let $e_{\boldsymbol k}$ count nonnegative excursions with that composition,
and set
\begin{equation}\label{eq:ambient}
 \E(\boldsymbol x)=\sum_{\boldsymbol k\ge0}
 e_{\boldsymbol k}\boldsymbol x^{\boldsymbol k}.
\end{equation}
The empty excursion contributes one. A nonzero coefficient requires
$\sum_i s_i k_i=0$, and
\begin{equation}\label{eq:diag}
 F_r(z)=\diag\E(z)
 :=\sum_{n\ge0}[x_1^n\cdots x_r^n]\E(\boldsymbol x)z^n.
\end{equation}
In particular, the dimension of this ambient generating function is fixed
when $r$ is fixed.

\subsection{The small-root formula, with a proof}
Temporarily mark total length by $t$, and write
\[
 P_{\boldsymbol x}(u)=\sum_{i=1}^r x_i u^{s_i}.
\]
The kernel equation is
\begin{equation}\label{eq:kernel}
 u^h\bigl(1-tP_{\boldsymbol x}(u)\bigr)=0.
\end{equation}
Of its $2h$ roots, exactly $h$ tend to zero as $t\to0$; these are called
small roots. This description may be understood in a Puiseux extension
with generic step weights.

\begin{lemma}[Kernel product]\label{lem:kernel}
If $u_1(t),\ldots,u_h(t)$ are the small roots, then
\begin{equation}\label{eq:rootproduct}
 \E(t\boldsymbol x)=
 \frac{(-1)^{h+1}}{t x_1}\prod_{j=1}^h u_j(t).
\end{equation}
Consequently $\E$ is an algebraic multivariate formal power series.
\end{lemma}
\begin{proof}
Let $G_j(t)$ count nonnegative paths ending at height $j$, and put
$G(t,u)=\sum_{j\ge0}G_j(t)u^j$. Appending one step gives
\[
 (1-tP_{\boldsymbol x}(u))G(t,u)
 =1-t\sum_{s_i<0}x_i\sum_{j=0}^{-s_i-1}G_j(t)u^{j+s_i}.
\]
After multiplication by $u^h$, the right-hand side is a monic polynomial
$N(t,u)$ of degree $h$. Substitution of any small root is legitimate in
the formal path equation and gives $N(t,u_j(t))=0$. Hence
$N(t,u)=\prod_{j=1}^h(u-u_j(t))$. Its constant coefficient is
$-t x_1G_0(t)$: only the most negative step, followed from height zero,
contributes to that coefficient. Equating constant coefficients proves
\eqref{eq:rootproduct}, since $G_0(t)=\E(t\boldsymbol x)$.

The roots are algebraic over $\Q(t,\boldsymbol x)$, so their product is
algebraic. An algebraic relation for $\E(t\boldsymbol x)$ can be specialized
to $t=1$ after removing any common factor $t-1$ from its coefficients.
The specialization is valid coefficientwise in the $\boldsymbol x$-adic
topology and yields a nonzero polynomial relation for $\E(\boldsymbol x)$.
\end{proof}
This is the classical directed-path kernel construction, included here to
make the composition marking explicit; its general development is due to
the directed-path literature, notably Banderier and Flajolet \cite{bf}.

\subsection{Why D-finiteness follows}
An algebraic multivariate formal series in characteristic zero is
D-finite: its derivatives belong to its finite algebraic function field,
so they span a finite-dimensional vector space over the rational function
field. Lipshitz's diagonal theorem then implies that the complete diagonal
\eqref{eq:diag} is D-finite \cite{lipshitz}.

Equivalently, for each fixed $r$ there exist polynomials
$q_0(n),\ldots,q_d(n)\in\Q[n]$, not all zero, with
\[
 \sum_{j=0}^d q_j(n)A_r(n+j)=0
\]
for all sufficiently large $n$. This proves the D-finiteness portion of
Theorem~\ref{thm:class}. Closure algorithms make this an effective existence
result, but they do not imply a small order, a small degree, or a fast
calculation of a useful recurrence.

\subsection{A concrete algebraic certificate for the fifth alphabet}
For steps $-2,-1,0,1,2$, give the weights the names $a,b,c,d,e$ respectively.
The ambient excursion series $X=\E(a,b,c,d,e)$ is the unique formal solution
with constant coefficient one of
\begin{align}\label{eq:sextic}
0={}&a^3e^3X^6+a^2e^2(c-1)X^5+(abde-a^2e^2)X^4\notag\\
 &+\bigl(ad^2+b^2e+2ae(1-c)\bigr)X^3
 +(bd-ae)X^2+(c-1)X+1.
\end{align}
Here the zero-step weight $c$ is a variable, not the OEIS asymptotic constant.

To verify the equation, factor the quartic kernel as
\[
 e(u^2-vu+q)(u^2-wu+q')
 =eu^4+du^3+(c-1)u^2+bu+a.
\]
The two small roots have product $q=-aX$, while
$q'=-1/(eX)$. Comparing the cubic and linear coefficients gives
\[
 v+w=-d/e,\qquad vq'+wq=-b/e.
\]
Solve these two equations for $v,w$, and substitute into
$q+q'+vw=(c-1)/e$. Clearing denominators gives \eqref{eq:sextic}.
This calculation is valid first for generic weights, and hence as a
polynomial identity. At $a=b=c=d=e=0$ the equation reduces to $1-X=0$,
with nonzero derivative in $X$; the formal implicit function theorem
therefore selects a unique branch. The fifth row is the complete diagonal
of this explicitly specified algebraic branch.

\section{The bridge logarithm and the algebraic constant}
\subsection{An exact identity retaining every composition}
Define the bridge generating function
\[
 B(t;\boldsymbol x)=\ct_u\frac{1}{1-tP_{\boldsymbol x}(u)}.
\]
Unlike excursions, bridges have no sign restriction on their intermediate
heights. Their composition coefficients are multinomial coefficients.

\begin{lemma}[Bridge logarithm]\label{lem:bridge}
As formal power series,
\begin{equation}\label{eq:bridgeidentity}
 B(t;\boldsymbol x)=1+t\frac{\mathrm d}{\mathrm dt}
             \log\E(t\boldsymbol x),
\end{equation}
and consequently
\begin{equation}\label{eq:logcoeff}
 \log\E(\boldsymbol x)=
 \sum_{\substack{\boldsymbol k\ne0\\\sum_i s_i k_i=0}}
 \frac{(|\boldsymbol k|-1)!}{\prod_i k_i!}
 \boldsymbol x^{\boldsymbol k}.
\end{equation}
In particular,
\begin{equation}\label{eq:exactreference}
 [x_1^n\cdots x_r^n]\log\E(\boldsymbol x)
 =\frac{1}{rn}\frac{(rn)!}{(n!)^r}\qquad(n\ge1).
\end{equation}
\end{lemma}
\begin{proof}
For small positive numerical weights and sufficiently small $t$, constant
term extraction is a contour integral. Its residues at the small kernel
roots give
\[
 B(t;\boldsymbol x)
 =\sum_{j=1}^h\frac{-1}{t u_jP'_{\boldsymbol x}(u_j)}
 =t\sum_{j=1}^h\frac{u_j'(t)}{u_j(t)}.
\]
The second equality follows by differentiating
$tP_{\boldsymbol x}(u_j(t))=1$. Taking the logarithmic derivative of
\eqref{eq:rootproduct} proves \eqref{eq:bridgeidentity}. The identity
extends formally to arbitrary weights.

Expanding $B$ in powers of $t$ and then expanding
$P_{\boldsymbol x}(u)^m$ gives the coefficient
$m!/\prod k_i!$ precisely when $|\boldsymbol k|=m$ and
$\sum s_i k_i=0$. Integrating $(B-1)/t$ supplies the factor $1/m$ and
proves \eqref{eq:logcoeff}. Its balanced specialization gives
\eqref{eq:exactreference}.
\end{proof}
The bridge--excursion relation is classical \cite{bf}. Its usefulness here
is that the diagonal of the logarithm is \emph{known exactly}. We can
calibrate a multivariate saddle calculation against that coefficient and
thereby avoid an unjustified guess about cyclic rotations or lattice
periods.

\subsection{The critical value}
Set $\boldsymbol p=(1/r,\ldots,1/r)$. In the probability normalization,
\[
 B(t;\boldsymbol p)=\sum_{m\ge0}\Prob(S_m=0)t^m.
\]
A finite-support, nondegenerate centered lattice walk has
$\Prob(S_m=0)=O(m^{-1/2})$. For completeness, this bound follows by Fourier
inversion: the characteristic function has quadratic modulus decay near
the finitely many points where its modulus is one, and has modulus bounded
strictly below one on their complement. The integral near each such point
is bounded by a Gaussian integral. Thus \eqref{eq:kappaseries} converges.
By monotone convergence in the nonnegative coefficient series,
\begin{equation}\label{eq:criticalvalue}
 \kappa_r=\lim_{t\uparrow1}\E(t\boldsymbol p)
 =\E(\boldsymbol p)>1.
\end{equation}

\begin{proposition}[Finite algebraic formula]\label{prop:kappa}
Let
\begin{equation}\label{eq:Q}
 Q_r(u)=\sum_{i=1}^r u^{s_i+h}-r u^h.
\end{equation}
The root $u=1$ has multiplicity exactly two. Apart from that root,
$Q_r$ has $h-1$ roots strictly inside the unit disk and $h-1$ strictly
outside, counting multiplicities. If the inside roots are
$\alpha_1,\ldots,\alpha_{h-1}$, then
\begin{equation}\label{eq:finiteconstant}
 \kappa_r=(-1)^{h+1}r\prod_{j=1}^{h-1}\alpha_j.
\end{equation}
An empty product is one.
\end{proposition}
\begin{proof}
At $u=1$, $P_r(1)=1$, $P'_r(1)=0$, and
$P''_r(1)=\sum_i p_i s_i^2>0$. Hence the multiplicity is exactly two.
If $|u|=1$ and $P_r(u)=1$, equality in the triangle inequality forces
$u^{s_i}=1$ for every $i$. The normalized step set has gcd one, so $u=1$.
There are no other roots on the unit circle.

For $0<t<1$, Rouch\'e's theorem applied on $|u|=1$ shows that
$u^h(1-tP_r(u))$ has exactly $h$ roots inside. Near $t=1$, its two roots
near one split as $1\pm c\sqrt{1-t}+O(1-t)$ with $c>0$; one is inside
and one outside. This leaves $h-1$ strictly inside in the limit. Formula
\eqref{eq:finiteconstant} is now the limit of \eqref{eq:rootproduct}.
It also proves algebraicity. Positivity was already obtained from
\eqref{eq:criticalvalue}.
\end{proof}
Internal multiple roots, should they occur, cause no difficulty: the
product of a separated cluster is analytic even when individual root
labels cannot be chosen analytically.

For the symmetric step sets in this paper one also has the real integral
\begin{equation}\label{eq:integralconstant}
 \log\kappa_r=-\frac1{2\pi}\int_0^{2\pi}
                       \log\bigl(1-P_r(e^{i\theta})\bigr)\dd\theta.
\end{equation}
Indeed, expand first with a factor $t<1$, integrate term by term, and let
$t\uparrow1$. The logarithmic singularity at $\theta=0$ is integrable.
The integrand is real because $P_r(e^{i\theta})$ is real and at most one.
Equation~\eqref{eq:integralconstant} is an alternative numerical expression,
not a replacement for the algebraic proof.

\section{Fixed-composition singularity analysis}\label{sec:saddle}
\subsection{Removing the two linear constraints}
Let $d=r-2$ and mark only the first $d$ step types in addition to total
length. Use logarithmic markers $\boldsymbol y=(y_1,\ldots,y_d)$ and set
\begin{align}
 P_{\boldsymbol y}(u)
 &=\sum_{i=1}^r p_i e^{y_i^*}u^{s_i},
 &y_i^*&=\begin{cases}y_i&i\le d,\\0&i>d.\end{cases}\label{eq:markedP}
\end{align}
If total length is $N$ and the first $d$ counts are known, the remaining
two counts are uniquely determined by length and total height zero,
because $s_{r-1}\ne s_r$. At $N=rn$, extracting $n$ occurrences of each
marked type therefore extracts the complete balanced composition.

The relevant coefficient functional is
\begin{equation}\label{eq:coefficientfunctional}
 \mathcal C_N(U)=[t^N z_1^n\cdots z_d^n]\,
 U\bigl(tp_1z_1,\ldots,tp_dz_d,tp_{r-1},tp_r\bigr),\quad N=rn.
\end{equation}
For $d=0$ the $z$ extraction is absent. Thus
\[
 \mathcal C_N(\E)=r^{-N}A_r(n),\qquad
 \mathcal C_N(\log\E)=r^{-N}\frac{(rn)!}{rn(n!)^r}.
\]

\subsection{The moving critical point and its Hessian}
The analytic implicit function theorem gives a unique $\tau(\boldsymbol y)$
near one solving
\begin{equation}\label{eq:criticalmoving}
 (u\partial_u)P_{\boldsymbol y}(u)\big|_{u=\tau(\boldsymbol y)}=0,
 \qquad \tau(0)=1.
\end{equation}
The required derivative in $\log u$ is the positive variance
\[
 V=\sum_i p_i s_i^2.
\]
Define
\begin{align}
 \rho(\boldsymbol y)&=P_{\boldsymbol y}(\tau(\boldsymbol y))^{-1},\notag\\
 \psi(\boldsymbol y)&=\log P_{\boldsymbol y}(\tau(\boldsymbol y)),\notag\\
 \Phi(\boldsymbol y)&=\psi(\boldsymbol y)-\sum_{a=1}^d p_a y_a.
 \label{eq:phase}
\end{align}
Then $\nabla\Phi(0)=0$. Its Hessian is
\begin{equation}\label{eq:H}
 H_{ab}=p_a\delta_{ab}-p_ap_b-
            \frac{p_as_a\,p_bs_b}{V}\qquad(1\le a,b\le d).
\end{equation}

\begin{lemma}[Nondegenerate composition saddle]\label{lem:hessian}
The matrix $H$ is positive definite, and, with
$\delta=s_r-s_{r-1}=2/g_r$,
\begin{equation}\label{eq:detH}
 \det H=\frac{\delta^2}{V}\prod_{i=1}^r p_i.
\end{equation}
For $d=0$, the determinant convention is $\det H=1$.
\end{lemma}
\begin{proof}
For numbers $v_1,\ldots,v_d$, extend $f_i=v_i$ for $i\le d$ by
$f_{r-1}=f_r=0$. Direct differentiation in \eqref{eq:criticalmoving} gives
\[
 v^{\mathsf T}Hv=\Var_p(f)-\frac{\Cov_p(f,s)^2}{V}.
\]
Cauchy--Schwarz makes this nonnegative. Equality would require $f_i=a+bs_i$
for every $i$. The two distinct steps with $f_i=0$ force $a=b=0$, hence
$v=0$.

For the determinant, let the random vector $X$ have components
$(\boldsymbol1_{\{i=1\}},\ldots,\boldsymbol1_{\{i=d\}},s_i)$ with
probability $p_i$. Its covariance matrix has Schur complement $H$ in the
last entry $V$, so its determinant is $V\det H$. A random vector supported
on $r$ affinely independent points in dimension $r-1$ has covariance
determinant equal to the product of its $r$ probabilities times the square
of the affine-simplex determinant. This identity follows by adjoining a
row of ones to the support matrix and applying the Schur complement to
its weighted Gram matrix. Here that affine determinant is, up to sign,
$s_r-s_{r-1}=\delta$. This proves \eqref{eq:detH}.
\end{proof}

\subsection{A uniform square-root expansion}
Put
\begin{equation}\label{eq:tiltedprob}
 q_i(\boldsymbol y)=
 \frac{p_i e^{y_i^*}\tau(\boldsymbol y)^{s_i}}
      {P_{\boldsymbol y}(\tau(\boldsymbol y))},\qquad
 V(\boldsymbol y)=\sum_i q_i(\boldsymbol y)s_i^2.
\end{equation}
The $q_i$ sum to one and have mean step zero. Multiplying weights by
$\tau^{s_i}$ does not change any excursion weight, so the critical value is
\[
 \kappa(\boldsymbol y)=\E(\boldsymbol q(\boldsymbol y)).
\]
It is analytic near $\boldsymbol y=0$ as a critical root-cluster product,
and $\kappa(0)=\kappa_r$.

Let $x=1-t/\rho(\boldsymbol y)$. The small root which reaches the critical
point has the expansion
\[
 u_0(t,\boldsymbol y)=\tau(\boldsymbol y)
 \left(1-c(\boldsymbol y)x^{1/2}+O(x)\right),\qquad
 c(\boldsymbol y)=\sqrt{\frac{2}{V(\boldsymbol y)}}.
\]
All the other small roots are a separated cluster. Their product, and the
factor $1/t$ in the kernel formula, are analytic at the critical point.
Therefore, uniformly for complex $\boldsymbol y$ sufficiently near zero,
\begin{align}
 \E(t;\boldsymbol y)
 &=\kappa(\boldsymbol y)
       \left(1-c(\boldsymbol y)x^{1/2}+O(x)\right),\label{eq:Elocal}\\
 \Lcal(t;\boldsymbol y):=\log\E(t;\boldsymbol y)
 &=\log\kappa(\boldsymbol y)-c(\boldsymbol y)x^{1/2}+O(x).
 \label{eq:Llocal}
\end{align}
Each expansion continues to arbitrary order in $x^{1/2}$. The notation
$\E(t;\boldsymbol y)$ means the marked specialization in
\eqref{eq:coefficientfunctional}, with $z_a=e^{y_a}$.

The square-root coefficients differ by exactly
$\kappa(\boldsymbol y)$. This analytic fact is the reason that the exact
logarithmic reference coefficient in \eqref{eq:exactreference} identifies
the fixed-composition leading constant.

\subsection{Why other phases do not contribute}
An asymptotic proof must control the full marker torus, not just a local
expansion. Here this control is particularly simple.

Suppose $|t|=|u|=|z_a|=1$ and a kernel root lies on $|u|=1$. Equality in
\[
 \left|t\sum_i p_i z_i^*u^{s_i}\right|\le1
\]
forces $t z_i^*u^{s_i}=1$ for every $i$. The last two unmarked types imply
$u^{\delta}=1$. All step differences are multiples of $\delta$, so then
$z_a=1$ for every marked type. If $r$ is odd, $\delta=1$ and $t=u=1$.
If $r$ is even, $\delta=2$ and the two possibilities are $(t,u)=(1,1)$
and $(-1,-1)$.

Thus the only marker saddle is $\boldsymbol z=\one$. The only length
period is
\begin{equation}\label{eq:period}
 \nu=\begin{cases}1&r\text{ odd},\\2&r\text{ even},\end{cases}
 \qquad\text{so that }\nu=\delta.
\end{equation}
In the even case every excursion has even length, and the two length
singularities give identical contributions on the admissible subsequence
$N=rn$.

More explicitly, on any compact part of the marker torus separated from
$\one$, no kernel root can cross $|u|=1$ for $|t|\le1$. By compactness
this remains true for $|t|<1+\epsilon$ with some uniform $\epsilon>0$.
The small-root product is analytic there, even if roots within that
cluster collide. It is nonzero away from $t=0$, since the kernel has
nonzero constant term for $t\ne0$. The removable value at $t=0$ is one.
Consequently both $\E$ and its logarithm have uniform analytic
continuations on that disk. Cauchy estimates make this part of the
coefficient integral exponentially smaller than the main term.

\subsection{Transfer and the Gaussian integral}
For a half-integer $\gamma$ which is not a nonnegative integer,
\begin{equation}\label{eq:transfer}
 [t^N](1-t/\rho)^\gamma
 =\rho^{-N}\frac{\Gamma(N-\gamma)}{\Gamma(-\gamma)\Gamma(N+1)}.
\end{equation}
The gamma quotient has a complete inverse-power expansion. Applying
\eqref{eq:transfer} to \eqref{eq:Elocal}--\eqref{eq:Llocal} gives local
leading amplitudes
\begin{equation}\label{eq:amplitudes}
 g_L(\boldsymbol y)=\frac{c(\boldsymbol y)}{2\sqrt\pi},\qquad
 g_E(\boldsymbol y)=\kappa(\boldsymbol y)g_L(\boldsymbol y),
\end{equation}
multiplying $\rho(\boldsymbol y)^{-N}N^{-3/2}$ and the period factor $\nu$.

For marker extraction put $\boldsymbol y=i\boldsymbol\theta$. The
exponential in the Cauchy integral is
\[
 \exp\bigl(N\Phi(i\boldsymbol\theta)\bigr),\qquad
 \Phi(i\boldsymbol\theta)=-\tfrac12\boldsymbol\theta^{\mathsf T}
 H\boldsymbol\theta+O(|\boldsymbol\theta|^3).
\]
The Hessian is positive definite. The leading Gaussian integral thus gives
\begin{align}
 \mathcal C_N(\Lcal)
 &\sim\frac{\nu g_L(0)}{(2\pi)^{d/2}\sqrt{\det H}}
                   N^{-(d+3)/2},\label{eq:Lcoefficient}\\
 \mathcal C_N(\E)
 &\sim\kappa_r\frac{\nu g_L(0)}{(2\pi)^{d/2}\sqrt{\det H}}
                   N^{-(d+3)/2}.
 \label{eq:Ecoefficient}
\end{align}
The ratio is therefore $\kappa_r$. Substituting the exact logarithmic
coefficient proves \eqref{eq:probability}.

There is also a useful independent normalization check. From
\eqref{eq:detH} and \eqref{eq:period},
\[
 \frac{\nu g_L(0)}{(2\pi)^{d/2}\sqrt{\det H}}
 =\frac{\nu}{\delta}\frac{1}{(2\pi)^{(r-1)/2}\sqrt{\prod_i p_i}}
 =\frac{r^{r/2}}{(2\pi)^{(r-1)/2}}.
\]
This is exactly the leading Stirling coefficient of
$r^{-N}(rn)!/(rn(n!)^r)$ when $N=rn$. In particular, no hidden period-two
factor is lost in the even alphabets.

\subsection{Remainders and all orders}\label{subsec:allorders}
We spell out why the preceding calculation proves more than a leading
heuristic. Fix any requested order $M$. The local kernel factorization
provides Puiseux expansions to an arbitrarily high finite order,
with analytic coefficient functions and uniform remainder estimates on
a common dented neighborhood of the moving singularity. This follows
from analytic factorization of the double root and analytic separation
of the other root clusters. The same neighborhoods work for sufficiently
small complex marker values.

After the length coefficient is extracted, Taylor-expand its analytic
amplitudes and the phase $\Phi$ to a finite order larger than that needed
for $N^{-M-1}$. In a neighborhood
$|\boldsymbol\theta|\le N^{-1/2+\varepsilon}$, substitute
$\boldsymbol\theta=\boldsymbol x/\sqrt N$. Taylor's theorem and the
positive-definite Gaussian majorant bound every discarded term by an
integrable polynomial times a Gaussian, with the specified inverse power
of $N$. One may first fix a sufficiently small $\varepsilon>0$ depending
on the requested truncation and then take sufficiently many Taylor
terms. In the remainder of the small marker neighborhood, the real part
of the phase is at most $-c|\boldsymbol\theta|^2$, giving an error smaller
than every power of $N^{-1}$. The compact complement was handled by the
analytic-continuation argument above.

Gaussian moments of odd total degree vanish. The half-integer powers in
the length expansion differ by integers after coefficient transfer.
Consequently only integer powers of $N^{-1}$ remain after the common
factor $N^{-(r+1)/2}$ is removed. This proves a full expansion with an
$O(N^{-M-1})$ relative remainder after $M$ corrections.

All finite Taylor jets involved are algebraic. Indeed, after replacing
$e^{y_a}$ by independent variables, the critical point solves polynomial
kernel and derivative equations with algebraic base values. Exponential
reparametrization has rational Taylor coefficients. Half-integer gamma
factors contribute one common $\sqrt\pi$, and Gaussian moments contribute
rational expressions in $H^{-1}$ times a common determinant factor.
After division by the leading coefficient, every correction is algebraic.

Finally, Stirling's expansion in
\eqref{eq:probability}, or directly \eqref{eq:Ecoefficient}, yields
\eqref{eq:full} and \eqref{eq:constant}. This proves
Theorem~\ref{thm:main}, including its all-order and algebraicity assertions.

\begin{remark}[What ``all orders'' means here]
Equation~\eqref{eq:full} determines the dominant exponential--power sector
with a proved finite-order remainder for every fixed truncation. It does
not assert convergence of the inverse-power series, identify every
exponentially small contribution, or establish Borel summability.
Those are separate questions.
\end{remark}

\section{Explicit leading constants}
The following values are obtained from the finite product, not from
fitting a growth curve:
\begin{center}
\begin{tabular}{crr}
\toprule
$r$ & $\kappa_r$ & $c(r)$ in \eqref{eq:oeisconj}\\
\midrule
2&2.000000000000000&1.000000000000000\\
3&3.000000000000000&0.866025403784439\\
4&1.384057356943304&0.244669085661435\\
5&1.909830056250526&0.213525491562421\\
6&1.239095083832872&0.089423985025306\\
7&1.583896616939167&0.074832081265314\\
8&1.173744890362129&0.036679527823817\\
\bottomrule
\end{tabular}
\end{center}
The printed decimals are numerical evaluations; the exact definition is
\eqref{eq:finiteconstant}.

For $r=2,3$, $h=1$, so the product is empty and $\kappa_2=2$,
$\kappa_3=3$. Directly,
\[
 A_2(n)=\frac{(2n)!}{n!(n+1)!},\qquad
 A_3(n)=\frac{(3n)!}{(n!)^3(n+1)}.
\]
For the second formula, remove the $n$ zero steps from a balanced Motzkin
excursion, obtaining a Dyck path with $n$ up and $n$ down steps; reinsert
the zero steps in $\binom{3n}{n}$ ways. These formulas give the checks
\[
 \alpha_{2,1}=-\frac98,\qquad \alpha_{3,1}=-\frac{11}{9}.
\]

For the fourth row,
\[
 Q_4(u)=(u-1)^2(u^4+2u^3+4u^2+2u+1).
\]
The quartic is reciprocal. On dividing by $u^2$ and putting $v=u+u^{-1}$,
it becomes $v^2+2v+2=0$. Pair the two roots strictly inside the unit disk
and multiply them. If $\varphi=(1+\sqrt5)/2$, the result is
\[
 \kappa_4=4(\varphi-\sqrt\varphi),\qquad
 c(4)=\frac{\varphi-\sqrt\varphi}{\sqrt2}.
\]
Substitution in \eqref{eq:full} reproduces the formula credited in the
OEIS entry to Kotesovec \cite{oeis4}.

For the fifth row,
\[
 Q_5(u)=(u-1)^2(u^2+3u+1).
\]
The unique strict inside root is $\alpha=(-3+\sqrt5)/2$. Hence
\[
 \kappa_5=\frac52(3-\sqrt5),\qquad c(5)=\frac{3\sqrt5-5}{8}.
\]
This is the leading constant already recorded in A215570 \cite{oeis5}.
For reproducibility, two larger reduced kernels are
\begin{align*}
 \frac{Q_6(u)}{(u-1)^2}
 &=u^8+2u^7+4u^6+6u^5+9u^4+6u^3+4u^2+2u+1,\\
 \frac{Q_7(u)}{(u-1)^2}
 &=u^4+3u^3+6u^2+3u+1.
\end{align*}
Together with the inside-root selection, these give exact algebraic
specifications of the constants for A215571 and A215593, whose checked
entries did not give asymptotic formulas \cite{oeis6,oeis7}.

\section{Why every nontrivial higher row is transcendental}\label{sec:transcendence}
The exponent and constant in Theorem~\ref{thm:main} imply a complete
algebraicity classification for the row generating functions.

\begin{proof}[Proof of the transcendence assertion in Theorem~\ref{thm:class}]
Set $R=r^{-r}$ and $G(x)=F_r(Rx)$. Its coefficients satisfy
\[
 [x^n]G(x)\sim K_r n^{-(r+1)/2},\qquad K_r>0.
\]
Assume for contradiction that $G$ is algebraic over $\Q(x)$. Every
ordinary derivative of $G$ is then algebraic and, near the algebraic
point $x=1$, has a Puiseux expansion with algebraic coefficients.

First let $r=2k+1\ge3$. Then
\[
 G^{(k)}(x)\sim K_r\log\frac1{1-x}\qquad(x\uparrow1).
\]
To justify this Abelian estimate, its coefficient sequence is asymptotic
to $K_r/n$ up to the harmless derivative index shift. Compare it, outside
any fixed finite initial segment, with $(K_r\pm\epsilon)/n$ and use
$\sum_{n\ge1}x^n/n=-\log(1-x)$. No Puiseux series can have a leading
logarithmic divergence. This is a contradiction.

Next let $r=2k\ge4$. Differentiation $k$ times gives coefficients
asymptotic to $K_r n^{-1/2}$. The same comparison argument and the
Gaussian integral, or a Riemann-sum estimate with $x=e^{-t}$, yield
\begin{equation}\label{eq:evenobstruction}
 G^{(k)}(x)\sim K_r\sqrt\pi\,(1-x)^{-1/2}.
\end{equation}
The leading Puiseux coefficient of an algebraic function over $\Q(x)$
at $x=1$ must be algebraic. But
\[
 K_r\sqrt\pi
 =\frac{\kappa_r}{\sqrt r\,2^{k-1/2}}\,\pi^{-(k-1)}
\]
is transcendental: its prefactor is nonzero algebraic, and $k-1$ is a
positive integer. If this number were algebraic, then $\pi^{k-1}$ and
hence $\pi$ would be algebraic, contradicting Lindemann's theorem
\cite{lindemann}. This contradicts \eqref{eq:evenobstruction}.

Finally $F_2(z)=(1-\sqrt{1-4z})/(2z)$ is algebraic, as expected.
\end{proof}
The argument is not a generic claim that all diagonals are transcendental.
It uses the particular positive leading asymptotic and its algebraic
constant. It also explains why the even rows require an arithmetic
argument in addition to their half-integer growth exponent.

\section{The first correction for the fifth row}\label{sec:correction}
\subsection{A differential formula for the ratio correction}
This section makes the all-order construction concrete. Its calculation
is independent of the conjectured A215570 recurrence.

Use the phase and marker notation from Section~\ref{sec:saddle}, and let
$C=H^{-1}$. Repeated indices in the formulas below are summed from $1$ to
$d$. Put
\[
 \ell(\boldsymbol y)=\log\kappa(\boldsymbol y),\qquad
 h_c(\boldsymbol y)=\log c(\boldsymbol y),\qquad
 T_{abc}=\partial_a\partial_b\partial_c\Phi(0).
\]
All derivatives without an argument are evaluated at zero.
For an analytic amplitude $g$, the order-$N^{-1}$ contribution in its
Gaussian saddle integral is
\[
 -\tfrac12 C_{ab}g_{ab}
 +\tfrac12 g_aT_{bcd}C_{ab}C_{cd}
 +g\,J,
\]
where $J$ is independent of the amplitude. This follows by expanding
$g(i\boldsymbol x/\sqrt N)$ and the cubic and quartic terms of
$N\Phi(i\boldsymbol x/\sqrt N)$, then using Gaussian moments. The signs
come from $i^2=-1$; in particular the second derivative term is negative.

Subtracting the normalized correction for $g_L$ from that for
$g_E=\kappa g_L$ cancels $J$ and gives
\begin{equation}\label{eq:saddlecorrection}
 \mathcal S=-\tfrac12 C_{ab}(\ell_{ab}+\ell_a\ell_b)
             -C_{ab}\ell_a(h_c)_b
             +\tfrac12\ell_aT_{bcd}C_{ab}C_{cd}.
\end{equation}
There is a second correction from the length singularity itself. Write
at $\boldsymbol y=0$
\begin{align*}
 \E(t)&=\kappa_r(1-cx^{1/2}+D x+E_3x^{3/2}+O(x^2)),\\
 \log\E(t)&=\log\kappa_r-cx^{1/2}+L_2x+L_3x^{3/2}+O(x^2),
 \qquad x=1-t.
\end{align*}
Then $D=L_2+c^2/2$. In the coefficient transfer formula the universal
$3/(8N)$ term cancels between the two functions, and
$\Gamma(-1/2)/\Gamma(-3/2)=-3/2$. It follows that the difference of their
relative length corrections is
\begin{equation}\label{eq:lengthcorrection}
 \mathcal B=-\frac32L_2-\frac{c^2}{4}.
\end{equation}
Combining the two extractions,
\begin{equation}\label{eq:ratiofirst}
 A_r(n)=\frac{\kappa_r}{rn}\frac{(rn)!}{(n!)^r}
 \left(1+\frac{\Delta_r}{rn}+O(n^{-2})\right),
 \qquad \Delta_r=\mathcal B+\mathcal S.
\end{equation}
Stirling's first correction then gives
\begin{equation}\label{eq:alphafirst}
 \alpha_{r,1}=\frac{1/r-r}{12}+\frac{\Delta_r}{r}.
\end{equation}
Equations~\eqref{eq:saddlecorrection}--\eqref{eq:alphafirst} are also an
implementable recipe for other fixed rows.

\subsection{Exact ingredients for \texorpdfstring{$r=5$}{r=5}}
Here $s=(-2,-1,0,1,2)$, $p_i=1/5$, $d=3$, and $V=2$. The Hessian and
inverse are
\[
 H=\begin{pmatrix}
 2/25&-2/25&-1/25\\
 -2/25&7/50&-1/25\\
 -1/25&-1/25&4/25
 \end{pmatrix},\qquad
 C=\begin{pmatrix}
 130&90&55\\90&70&40\\55&40&30
 \end{pmatrix}.
\]
The large entries in $C$ are a consequence of marking the first three
letters while reconstructing the final two. They are not evidence of a
degenerate saddle: $H$ is strictly positive definite.

For compact formulas define the centered residual markers
\begin{equation}\label{eq:residualmarkers}
 R_a(i)=\boldsymbol1_{\{i=a\}}-p_a-\frac{p_as_a}{V}s_i.
\end{equation}
Then
\begin{equation}\label{eq:thirdandh}
 T_{abc}=\sum_i p_iR_a(i)R_b(i)R_c(i),\qquad
 (h_c)_a=-\frac{1}{2V}\sum_i p_i s_i^2R_a(i).
\end{equation}
The cubic identity is the third derivative of the profiled log moment
generating function in \eqref{eq:phase}; derivatives of the optimized
height coordinate cancel because each residual marker is orthogonal to
the height step.

There is an especially simple formula for $\kappa(\boldsymbol q)$ for
nearby centered five-step probabilities. Write
\[
 A=q_{-2},\quad E=q_{2},\quad B=q_1+2q_2=q_{-1}+2q_{-2}.
\]
The equality for $B$ is the zero-drift condition. The critical kernel
factors as
\[
 \sum_{s=-2}^2q_su^{s+2}-u^2=(u-1)^2(Eu^2+Bu+A),
\]
so
\begin{equation}\label{eq:kappafivegeneral}
 \kappa(\boldsymbol q)=\frac{2}{B+\sqrt{B^2-4AE}}.
\end{equation}
The square root is positive at the uniform probabilities and continued
analytically nearby.

For transparency, let $v=\log\tau$. Differentiating the zero-drift equation
and \eqref{eq:tiltedprob} gives
\begin{align}
 v_a&=-p_as_a/V, &
 v_{ab}&=-V^{-1}\sum_i p_i s_iR_a(i)R_b(i),\notag\\
 (q_i)_a&=p_iR_a(i),&
 (q_i)_{ab}&=p_i\bigl(R_a(i)R_b(i)+s_i v_{ab}-H_{ab}\bigr).
 \label{eq:qjets}
\end{align}
Equations~\eqref{eq:kappafivegeneral}--\eqref{eq:qjets} determine the needed
derivatives using only rational arithmetic and $\sqrt5$. Their values are
\[
 (\ell_a)=\left(-\frac3{10}+\frac{\sqrt5}{10},\quad
                   \frac15-\frac{\sqrt5}{10},\quad\frac15\right),
 \qquad ((h_c)_a)=\left(-\frac1{10},\frac1{20},\frac1{10}\right),
\]
and
\[
 (\ell_{ab})=\begin{pmatrix}
 \frac2{25}-\frac{4\sqrt5}{125}&-\frac2{25}+\frac{11\sqrt5}{250}&-\frac1{25}\\[2pt]
 -\frac2{25}+\frac{11\sqrt5}{250}&\frac7{50}-\frac{17\sqrt5}{250}&-\frac1{25}\\[2pt]
 -\frac1{25}&-\frac1{25}&\frac4{25}
 \end{pmatrix}.
\]
Substitution in \eqref{eq:saddlecorrection}, using \eqref{eq:thirdandh},
now gives the exact contraction
\begin{equation}\label{eq:Sfive}
 \mathcal S=-\frac{11}{4}+\sqrt5.
\end{equation}

To obtain $L_2$, note that the uniform law is symmetric. The logarithm of
the critical small root has no $x$ term: in the coordinate $u=e^w$,
$P(e^w)$ is even in $w$, so $w=-c\sqrt{x}+O(x^{3/2})$. The other small
root is $\alpha=(-3+\sqrt5)/2$. Differentiating its kernel equation yields
\[
 L_2=1+\frac{1}{\alpha P'_5(\alpha)}=1-\frac{\sqrt5}{5}.
\]
Since $c^2=2/V=1$, equation~\eqref{eq:lengthcorrection} gives
\[
 \mathcal B=-\frac74+\frac{3\sqrt5}{10},\qquad
 \Delta_5=-\frac92+\frac{13\sqrt5}{10}.
\]
Finally \eqref{eq:alphafirst} gives
\[
 \alpha_{5,1}=-\frac{13}{10}+\frac{13\sqrt5}{50}
             =-\frac{13(5-\sqrt5)}{50}.
\]
This proves \eqref{eq:fivecorrection}.

All higher finite jets at the uniform point stay in $\Q(\sqrt5)$:
the critical root-cluster data have that field, $c=1$, and all implicit
recursions have nonzero denominators in the same field. Gaussian and
Stirling corrections introduce only rational operations. This proves the
field assertion in Theorem~\ref{thm:five}.

\section{A fixed-composition contact law}
The same saddle comparison provides distributional information, not only
total enumeration. An \emph{arch} is a nonempty excursion with no interior
return to height zero. A zero step at height zero is itself an arch.
Every nonempty excursion decomposes uniquely into a sequence of arches.
Hence the arch generating function is
\[
 H_0(\boldsymbol x)=1-\E(\boldsymbol x)^{-1}.
\]
Let $J_n$ be the number of arches of a uniformly chosen balanced
$r$-letter excursion, and let $C_n=J_n-1$ be its number of interior
returns to zero.

\begin{theorem}[Contact limit under exact composition]\label{thm:contacts}
For each fixed $r\ge2$ and $k\ge0$,
\begin{equation}\label{eq:contactlaw}
 \Prob(C_n=k)\longrightarrow
 \frac{k+1}{\kappa_r^2}\left(1-\frac1{\kappa_r}\right)^k.
\end{equation}
The probability generating functions converge in a complex neighborhood
of $v=1$, and
\begin{equation}\label{eq:meancontacts}
 \Exp C_n\longrightarrow 2(\kappa_r-1).
\end{equation}
\end{theorem}
\begin{proof}
For any analytic function $f$ near $\kappa_r$, the leading square-root
coefficient of $f(\E)$ is $f'(\kappa_r)$ times that of $\E$. The same
coefficient-extraction proof therefore gives the leading ratio
$f'(\kappa_r)$, provided the composition is analytic on the necessary
neighborhoods. For fixed $k$, apply this with
$f(E)=(1-E^{-1})^{k+1}$. It counts exactly $k+1$ arches, and the ratio is
\[
 f'(\kappa_r)=\frac{k+1}{\kappa_r^2}
                    (1-\kappa_r^{-1})^k.
\]
These nonnegative limiting probabilities sum to one.

For the stronger generating-function statement, mark arches by $v$:
\[
 \E_v=\frac{1}{1-vH_0}.
\]
At the critical weights, $H_0=\eta:=1-\kappa_r^{-1}<1$. On the closed
coefficient polydisk, nonnegative coefficients give $|H_0|\le\eta$.
For $|v|<1/\eta$ the denominator is consequently bounded away from zero;
this remains true in a suitable local continuation. The preceding
analysis is therefore uniform for $v$ near one. Its derivative ratio is
\[
 \Exp(v^{J_n})\longrightarrow
 \frac{v}{\kappa_r^2(1-\eta v)^2}.
\]
Divide by $v$ and differentiate at one to obtain \eqref{eq:meancontacts}.
\end{proof}
A related negative-binomial contact law is classical when only length is
conditioned \cite{bf}. The statement here imposes the complete step
composition. Its proof explains why the same critical value survives this
additional conditioning.

\section{Inverting the asymptotic expansion}\label{sec:inverse}
\subsection{An exact Lambert core}
For fixed $r$ define
\[
 \lambda=r\log r,\qquad \beta=\frac{r+1}{2},\qquad K=K_r.
\]
Then the proved expansion has the form
\begin{equation}\label{eq:genericexpansion}
 A_r(n)\sim K e^{\lambda n}n^{-\beta}
          \left(1+\frac{\alpha_1}{n}+\frac{\alpha_2}{n^2}+\cdots\right).
\end{equation}
The large real solution of the \emph{leading model}
$Y=Ke^{\lambda x}x^{-\beta}$ is exactly
\begin{equation}\label{eq:lambertcore}
 x_0(Y)=-\frac{\beta}{\lambda}
 W_{-1}\!\left(-\frac{\lambda}{\beta}
                  \left(\frac K Y\right)^{1/\beta}\right).
\end{equation}
The lower real branch $W_{-1}$ is essential: it selects $x_0\to\infty$.
The argument lies in $(-e^{-1},0)$ for all sufficiently large $Y$.
This exact-core organization parallels the inversion framework inspected
in ProveIt \cite{proveit}; no unverified asymptotic interpolation is needed.

\subsection{Corrections and a remainder statement}
Write the formal logarithm of the unit correction as
\[
 \log\left(1+\sum_{j\ge1}\alpha_jx^{-j}\right)
       =\sum_{j\ge1}b_jx^{-j},\qquad
 b_1=\alpha_1,\quad b_2=\alpha_2-\tfrac12\alpha_1^2.
\]
Seek $x=x_0+\sum_{j\ge1}d_jx_0^{-j}$ in
\begin{equation}\label{eq:inversecorrectioneq}
 \lambda(x-x_0)-\beta\log(x/x_0)+\sum_{j\ge1}b_jx^{-j}=0.
\end{equation}
At every order the new coefficient occurs with nonzero multiplier
$\lambda$, so the recursion is triangular. Its first terms are
\begin{equation}\label{eq:inversefirst}
 x=x_0-\frac{b_1}{\lambda x_0}
 -\frac{b_2/\lambda+\beta b_1/\lambda^2}{x_0^2}
 +O(x_0^{-3}).
\end{equation}

This has a precise interpretation on the original sequence. Given any
fixed order, evaluate the truncated inverse at $Y=A_r(n)$. The forward
expansion has a controlled logarithmic residual, while the derivative of
$\lambda x-\beta\log x$ is $\lambda-\beta/x$, bounded below by
$\lambda/2$ for large $x$. The mean value theorem transports the residual
to an error of the same inverse-power order in the recovered $n$.
Thus the inversion is a theorem about the sequence values and their
asymptotic inverse, rather than a claim that a canonical analytic
interpolation of the integers has been constructed.

For the fifth row, \eqref{eq:inversefirst} begins
\begin{equation}\label{eq:inversefive}
 n=x_0+\frac{13(5-\sqrt5)}{250\log5}\frac1{x_0}
                  +O(x_0^{-2}),\qquad Y=A_5(n),
\end{equation}
where $\beta=3$, $\lambda=5\log5$, and
$K=(3\sqrt5-5)/(8\pi^2)$ in \eqref{eq:lambertcore}.

\subsection{Integer thresholds are a separate operation}
For arbitrary real $Y$, define
\[
 T_r(Y)=\min\{n\ge0:A_r(n)\ge Y\}.
\]
Appending a fixed balanced excursion block containing one of each step
is an injection from the $n$th class into the $(n+1)$st. Hence $A_r(n)$
is nondecreasing; its proved asymptotic also gives
$A_r(n+1)/A_r(n)\to r^r>1$, so it is eventually strictly increasing.

The continuous inverse predicts $T_r(Y)$ to bounded additive error.
Taking a ceiling gives the exact threshold when an explicit inverse-error
bound separates the predicted value from the relevant integers. Without
that separation there is no universally valid exact rounding rule: $Y$
can be arbitrarily close to a jump value $A_r(n)$. This is the discrete
staircase issue emphasized in the repository's inversion documentation,
and it must not be hidden in an $o(1)$ statement about a continuous model.

\section{Unequal but centered compositions}
The main construction is not restricted to equal multiplicities. Keep the
step set \eqref{eq:steps}, and let $p_1,\ldots,p_r$ be fixed positive
rational numbers with
\[
 \sum_i p_i=1,\qquad \sum_i p_i s_i=0.
\]
For admissible $N$, let $e(N\boldsymbol p)$ count excursions with exactly
$Np_i$ occurrences of step $i$. Define
$\kappa(\boldsymbol p)=\E(\boldsymbol p)$ by the same critical kernel
product, now using these probabilities.

\begin{corollary}[Centered rational rays]\label{cor:rays}
Along integers $N$ for which all $Np_i$ are integral,
\begin{equation}\label{eq:ray}
 e(N\boldsymbol p)\sim
 \frac{\kappa(\boldsymbol p)}{N}\frac{N!}{\prod_i(Np_i)!}.
\end{equation}
There is a full inverse-power expansion in $N$ with algebraic normalized
coefficients. The coefficient series on every fixed positive integral
ray is D-finite.
\end{corollary}
\begin{proof}
Replace the uniform probabilities everywhere in Section~\ref{sec:saddle}
by the specified $p_i$. Centering keeps the critical height coordinate at
one. The Hessian proof, its determinant, and the marker-torus equality
argument are unchanged. The exact coefficient of $\log\E$ is again the
multinomial divided by $N$. Its ratio with the excursion coefficient gives
\eqref{eq:ray}, and the finite-jet argument gives all orders.

For the final statement, a fixed rational-ray extraction can be implemented
by coefficient sections and a diagonal. Sections are finite root-of-unity
filters after a monomial substitution; algebraicity is preserved by these
operations, and the subsequent diagonal is D-finite. Equivalently one
may use the multivariate D-finite residue closure for the fixed linear
coefficient constraints. The dimensions and integer ray vector remain
fixed.
\end{proof}
This extension concerns excursions relative to the \emph{same} height
functional $s_i$. Centering ensures the final height is zero. Compositions
with nonzero mean would describe a different endpoint problem and are
not covered by this statement.

\section{Reproducible checks and their scope}\label{sec:checks}
\subsection{Exact finite enumeration}
The accompanying program \path{code/verify.py} uses a direct composition
recurrence. For a count vector $\boldsymbol k$, let $D(\boldsymbol k)$
count nonnegative paths with that composition, regardless of final height.
Then
\[
 D(\boldsymbol0)=1,\qquad
 D(\boldsymbol k)=0\ \text{if }\sum_i s_i k_i<0,
\]
and for every other nonempty vector,
\[
 D(\boldsymbol k)=\sum_{i:k_i>0}D(\boldsymbol k-\boldsymbol e_i).
\]
The recurrence follows by deleting the last step. Its implementation stores
a finite rectangular count-vector box in mixed-radix order, so every
predecessor is already available. The diagonal gives $A_r(n)$.

The executed ranges were
\begin{center}
\begin{tabular}{ccc}
\toprule
$r$&Independently generated $n$&Rectangular states\\
\midrule
2&0--24&625\\
3&0--24&15,625\\
4&0--18&130,321\\
5&0--16&1,419,857\\
6&0--6&117,649\\
7&0--4&78,125\\
\bottomrule
\end{tabular}
\end{center}
The first two rows were compared with their factorial formulas. The other
rows were compared with the available initial values in the corresponding
OEIS entries. All comparisons passed.

An independent exact check of the bridge logarithm uses the identity
$D_{\rm Euler}\E=\E D_{\rm Euler}\log\E$. Coefficientwise this is
\begin{equation}\label{eq:expcheck}
 |\boldsymbol k|e_{\boldsymbol k}
 =\sum_{\substack{0<\boldsymbol j\le\boldsymbol k\\s\cdot\boldsymbol j=0}}
 \frac{|\boldsymbol j|!}{\prod_i j_i!}\,
 e_{\boldsymbol k-\boldsymbol j}.
\end{equation}
Exact rational arithmetic in this identity agreed with the direct path
recurrence for 120 nonempty balanced vectors across five test boxes.
A separate height-based weighted enumeration verified the sextic
\eqref{eq:sextic} through length 24 at weights $(2,3,5,7,11)$.
The polynomial elimination proving that sextic is also reproduced in
\path{code/derive_alpha5.py}.

\subsection{Asymptotic and inverse diagnostics}
For $A_5$, write
$L_5(n)=K_5 5^{5n}n^{-3}$ and $R_n=A_5(n)/L_5(n)$. The following values
illustrate the independently derived first correction:
\begin{center}
\begin{tabular}{rrrr}
\toprule
$n$&$R_n$&$n(R_n-1)$&$n^2(R_n-1-\alpha_{5,1}/n)$\\
\midrule
4&0.838539713679&-0.645841145284&0.291124722264\\
8&0.914884186065&-0.680926511482&0.301566514946\\
12&0.942234769960&-0.693182760477&0.305274784481\\
16&0.956285986890&-0.699424209762&0.307169857409\\
20&0.964839683966&-0.703206320687&0.308320103255\\
50&0.985752002111&-0.712399894449&0.311121570077\\
\bottomrule
\end{tabular}
\end{center}
The predicted limit of the middle diagnostic is
\[
 \alpha_{5,1}=-0.7186223258500546789\ldots.
\]
The values for $n\le16$ were generated independently by the supplied
program. The $n=20,50$ values were read from the published OEIS b-file;
they were \emph{not} independently enumerated here \cite{bfile5}.
The last column is a numerical check of a second-order residual, not a
claimed formula or a numerical certificate for $\alpha_{5,2}$.

At $Y=A_5(n)$, the Lambert core and its first correction give
\begin{center}
\begin{tabular}{rrr}
\toprule
$n$&$x_0(Y)-n$&$x_0(Y)-\alpha_{5,1}/(5\log5\,x_0(Y))-n$\\
\midrule
8&$-1.1595249025\cdot10^{-2}$&$-4.1641746283\cdot10^{-4}$\\
16&$-5.6870504107\cdot10^{-3}$&$-1.0375143173\cdot10^{-4}$\\
50&$-1.7966866816\cdot10^{-3}$&$-1.0601877690\cdot10^{-5}$\\
\bottomrule
\end{tabular}
\end{center}
The observed scales are consistent with the proved first- and second-order
inverse errors. Numerical agreement is a useful implementation check,
not a substitute for the asymptotic proof.

\subsection{What the computations do not certify}
The package does not verify a Lean formalization, prove the guessed
A215570 recurrence, find the minimal recurrence for any higher row, or
supply interval-certified decimal enclosures for every kernel root.
Its exact tests validate finite combinatorial and algebraic calculations.
Its root computations and asymptotic tables are high-precision numerical
diagnostics. The infinite statements are supported by the mathematical
proofs above.

\section{Further research questions and topics}\label{sec:future}
\subsection{Certify the explicit A215570 recurrence}
The sextic \eqref{eq:sextic}, its specified formal branch, and a diagonal
representation give a concrete input for creative telescoping. The next
objective is to compute a certified annihilator and prove that the
order-three operator in Kauers--Koutschan Conjecture 15 is a right factor
annihilating the same solution, with all initial conditions and singular
indices checked. Merely matching additional terms would not accomplish
this. A useful outcome would be a compact certificate independently
checkable without a large symbolic-computation environment.

\subsection{Find and explain the fourth-row recurrence}
A215562 is D-finite by the proof here, but its minimal recurrence and its
computational complexity remain to be determined. The even-step parity
and the larger kernel degree should be incorporated before elimination,
not left for a generic telescoper to rediscover. What upper bounds on
recurrence order and coefficient degree can be proved as functions of $r$?
Do the even and odd alphabets have qualitatively different minimal orders?

\subsection{Automate all-order coefficient extraction}
The proof supplies finite algebraic operations, but a practical system
should exploit residual-marker coordinates, symmetry, and Wick
contractions to avoid dense high-dimensional tensors. An immediate target
is exact $\alpha_{r,1},\alpha_{r,2},\alpha_{r,3}$ for $r=4,6,7,8$, together
with machine-readable algebraic-number specifications. For $r=5$, the
field $\Q(\sqrt5)$ makes an especially clean regression suite.

\subsection{Determine the exponentially small sectors}
The dominant expansion does not identify subdominant critical points,
connection constants, or Stokes phenomena. Once an explicit differential
or recurrence operator is available, compare its formal solutions with
complex kernel saddles. Which exponentially small contributions are
actually present in the combinatorial solution? Can optimal truncation
and rigorous exponentially improved error bounds be obtained without
first computing a minimal recurrence?

\subsection{Let the alphabet grow}
Our estimates keep $r$ fixed. When $r=r(n)$ grows, kernel-root separation,
Gaussian dimension, and constants in the remainder estimates all change.
Determine useful uniform ranges, and analyze the large-$r$ behavior of
$\kappa_r$ separately along even and odd alphabets. This would connect the
fixed-row problem to the less regular column direction of A215561.

\subsection{Approach the boundary of the composition simplex}
In Corollary~\ref{cor:rays}, positivity of all probabilities is essential
for the full-dimensional saddle. If some $p_i$ tend to zero, the determinant
in \eqref{eq:detH} degenerates and the effective alphabet may change.
Find uniform transition formulas interpolating between adjacent alphabet
sizes and identify the scaling regimes in which a different asymptotic
power appears.

\subsection{Refine path statistics under exact composition}
The contact law is only one consequence of the marked kernel. Natural
next targets are its finite-$n$ corrections, the maximum height, the area
under the excursion, and functional convergence after diffusive scaling.
These require additional arguments: neither the count asymptotic nor the
contact law alone proves a Brownian-excursion limit or independence of
other statistics.

\subsection{Formalize the exact-to-analytic interface}
A staged ProveIt development could first formalize the word--walk
bijection, the finite-state recurrences, the bridge coefficient identity,
and the formal algebraic kernel relation. A second stage would address
algebraic-series diagonal closure and critical-root selection. The most
substantial analytic stage is a reusable uniform marked square-root
transfer theorem with a nondegenerate composition saddle and explicit
lattice periods. Keeping these stages separate would prevent finite
computations from being mistaken for proofs of analytic hypotheses.

\section{Conclusion}
The fixed-row problem is governed by a finite one-dimensional step set,
not by an intrinsically growing transfer matrix. Marking composition
before taking a diagonal makes this structure visible. The resulting
algebraic ambient series gives D-finiteness; its bridge logarithm gives
an exact reference coefficient; and a marked critical-point calculation
converts that coefficient into the conjectured row asymptotic with an
algebraic constant.

The same analysis yields a complete dominant inverse-power expansion,
a concrete correction for A215570, a contact distribution, and an
asymptotic inversion procedure. A separate arithmetic argument shows
that every row generating function above the Catalan case is
transcendental. These results settle the stated fixed-row growth
conjecture and the recurrence-existence issue while leaving the explicit
recurrence certificates and exponentially small sectors as sharply
specified next problems.

\appendix
\section{A reproducible all-order construction}\label{app:algorithm}
For clarity, the finite procedure implicit in the proof can be organized
as follows. The input is a fixed $r$ and a requested relative order $M$.

\paragraph{Critical algebraic data.}
Construct $Q_r$, remove $(u-1)^2$, and isolate the strict inside cluster.
Represent its symmetric functions in an algebraic number field. Solve
\eqref{eq:criticalmoving} recursively for a finite Taylor jet of
$\tau(\boldsymbol y)$; obtain $\rho$, $\Phi$, and the tilted probabilities.
Internal root collisions can be handled by the coefficients of the
cluster factor instead of individual root labels.

\paragraph{Puiseux data.}
Set $t=\rho(\boldsymbol y)(1-w^2)$. Substitute
$u_0=\tau(\boldsymbol y)(1+\sum_{j\ge1}a_j(\boldsymbol y)w^j)$ in the
kernel, choosing $a_1=-\sqrt{2/V(\boldsymbol y)}$. The remaining
coefficients are obtained recursively. Multiply by the analytic
noncritical cluster product and the factor $1/t$ to form the needed
Puiseux jet of $\E$ and, independently, of $\log\E$.

\paragraph{Length coefficients.}
For every odd power $w^{2j+1}$, use \eqref{eq:transfer} and the gamma-ratio
expansion. The even powers belong to the local analytic part. Retain
enough terms for the requested order, and include the length-period
factor $\nu$.

\paragraph{Composition coefficients.}
Expand the phase about zero, substitute
$\boldsymbol y=i\boldsymbol x/\sqrt N$, and multiply the amplitude jets.
Replace each even monomial by its centered Gaussian moment with covariance
$H^{-1}$; discard odd monomials. Wick's rule evaluates these moments as
finite sums over pairings. The determinant and common $\pi$ factor are
kept outside the normalized series.

\paragraph{Calibration and inversion.}
Divide the resulting excursion expansion by the logarithmic reference
expansion if desired. The exact coefficient \eqref{eq:exactreference}
and Stirling's series are an independent normalization check. Convert
$N=rn$, compute the formal logarithm of the unit series, and use
\eqref{eq:inversecorrectioneq} for inverse coefficients.

Every stage is finite at fixed $r,M$. The procedure proves computability,
not an optimal complexity bound or a claim that a complete general-purpose
implementation is included in this package. The supplied exact symbolic
implementation covers the first correction for $r=5$.

\section{Package contents and build instructions}
The archive contains this editable source and its compiled PDF, the two
Python programs, their saved outputs, and source/provenance notes.
From the package directory, run
\begin{quote}\ttfamily
python -m pip install -r requirements.txt\\
python code/verify.py\\
python code/derive\_alpha5.py\\
pdflatex -interaction=nonstopmode -halt-on-error article.tex\\
pdflatex -interaction=nonstopmode -halt-on-error article.tex
\end{quote}
The computations use arbitrary-precision integers, rational arithmetic,
SymPy, and mpmath. They need no network access. Source URLs and the origin
of the two externally supplied large-index values are recorded in
\path{SOURCES.md}. No external paper is redistributed in the archive.

\begin{thebibliography}{99}
\bibitem{oeisarray}
The OEIS Foundation Inc., \emph{A215561}, array of prefix-average multiset
permutations. Entry by Alois P. Heinz; fixed-row asymptotic conjecture by
Vaclav Kotesovec, September 7, 2016.
\url{https://oeis.org/A215561}. Checked October 1, 2026.

\bibitem{oeis4}
The OEIS Foundation Inc., \emph{A215562}. Leading asymptotic formula by
Vaclav Kotesovec, January 31, 2015.
\url{https://oeis.org/A215562}. Checked October 1, 2026.

\bibitem{oeis5}
The OEIS Foundation Inc., \emph{A215570}. Leading asymptotic formula by
Vaclav Kotesovec, September 6, 2016; recurrence conjecture by Manuel
Kauers and Christoph Koutschan, March 2, 2023.
\url{https://oeis.org/A215570}. Checked October 1, 2026.

\bibitem{oeis6}
The OEIS Foundation Inc., \emph{A215571}.
\url{https://oeis.org/A215571}. Checked October 1, 2026.

\bibitem{oeis7}
The OEIS Foundation Inc., \emph{A215593}.
\url{https://oeis.org/A215593}. Checked October 1, 2026.

\bibitem{bfile5}
Manuel Kauers and Christoph Koutschan, with terms 0--47 credited in the
OEIS entry to Vaclav Kotesovec, \emph{Table of A215570, $0\le n\le50$}.
\url{https://oeis.org/A215570/b215570.txt}. Checked October 1, 2026.

\bibitem{kk}
M. Kauers and C. Koutschan,
\emph{Some D-finite and some possibly D-finite sequences in the OEIS},
Journal of Integer Sequences \textbf{26} (2023), Article 23.4.5,
especially Section 6.1 and Conjecture 15, pp.~29--30.
\url{https://cs.uwaterloo.ca/journals/JIS/VOL26/Koutschan/kout4.pdf}.
Also arXiv:2303.02793.

\bibitem{bf}
C. Banderier and P. Flajolet,
\emph{Basic analytic combinatorics of directed lattice paths},
Theoretical Computer Science \textbf{281} (2002), 37--80.
DOI: \href{https://doi.org/10.1016/S0304-3975(02)00007-5}{10.1016/S0304-3975(02)00007-5}.
\url{https://www-lipn.univ-paris13.fr/~banderier/Papers/tcs_banderier_flajolet_2002.pdf}.

\bibitem{lipshitz}
L. Lipshitz,
\emph{The diagonal of a D-finite power series is D-finite},
Journal of Algebra \textbf{113} (1988), 373--378.
DOI: \href{https://doi.org/10.1016/0021-8693(88)90166-4}{10.1016/0021-8693(88)90166-4}.

\bibitem{lindemann}
F. Lindemann, \emph{Ueber die Zahl $\pi$},
Mathematische Annalen \textbf{20} (1882), 213--225.
DOI: \href{https://doi.org/10.1007/BF01446522}{10.1007/BF01446522}.

\bibitem{proveit}
V. Reshetnikov, \emph{ProveIt}. Repository snapshot:\\
{\small\texttt{40b17f2fc79a10a6fc1a50ca9c0c3d6f77be69e5}}.\par
Inspected source: the transseries-and-inversion README, at the path
identified in Section~1.3.
\url{https://github.com/VladimirReshetnikov/ProveIt/tree/40b17f2fc79a10a6fc1a50ca9c0c3d6f77be69e5/Analysis/Transseries/docs/series-and-transseries/Transseries_And_Inversion}.
Accessed through the GitHub connector on October 1, 2026.
\end{thebibliography}
\end{document}
